%==============================================================
%  NOTATION — Macro toán học dùng chung cho toàn báo cáo.
%==============================================================

%--------------------------------------------------------------
% 0. KHAI BÁO CÁC GÓI PHỤ THUỘC (Đặt ở main.tex nếu chưa có)
%--------------------------------------------------------------
% \usepackage{amsmath, amssymb, amsfonts}
% \usepackage{bm}
% \usepackage{mathtools}

%--------------------------------------------------------------
% 1. TẬP HỢP SỐ VÀ MIỀN
%--------------------------------------------------------------
\newcommand{\R}{\mathbb{R}}
\newcommand{\N}{\mathbb{N}}
\newcommand{\Z}{\mathbb{Z}}
\newcommand{\E}{\mathbb{E}}                % Kỳ vọng: \E[\theta_k^2]
\newcommand{\Pprob}{\mathbb{P}}            % Xác suất
\newcommand{\dom}{\Omega}                  % Miền không gian
\newcommand{\bdry}{\partial\Omega}         % Biên của miền không gian
\newcommand{\timedom}{[0,T]}              % Miền thời gian

%--------------------------------------------------------------
% 2. VECTOR VÀ MA TRẬN (in đậm)
%--------------------------------------------------------------
% Chữ cái thường (Vector)
\newcommand{\ba}{\bm{a}}
\newcommand{\bb}{\bm{b}}
\newcommand{\bc}{\bm{c}}
\newcommand{\bd}{\bm{d}}
\newcommand{\be}{\bm{e}}
\newcommand{\bg}{\bm{g}}                   % Gradient
\newcommand{\bh}{\bm{h}}                   % Trạng thái ẩn (hidden state)
\newcommand{\bp}{\bm{p}}                   % Hướng tìm kiếm
\newcommand{\br}{\bm{r}}                   % Residue / Điểm dữ liệu
\newcommand{\bs}{\bm{s}}
\newcommand{\bu}{\bm{u}}                   % Nghiệm / Biến phụ thuộc
\newcommand{\bv}{\bm{v}}
\newcommand{\bw}{\bm{w}}
\newcommand{\bx}{\bm{x}}                   % Toạ độ không gian
\newcommand{\by}{\bm{y}}                   % Đầu ra / Toạ độ
\newcommand{\bz}{\bm{z}}                   % Giá trị kích hoạt (pre-activation)

% Chữ cái hoa (Ma trận)
\newcommand{\bA}{\bm{A}}
\newcommand{\bB}{\bm{B}}
\newcommand{\bD}{\bm{D}}                   % D^{(l)} = diag(sigma'(z^{(l)}))
\newcommand{\bF}{\bm{F}}                   % Ma trận Fisher Information
\newcommand{\bG}{\bm{G}}                   % Truy hồi Jacobi theo biến đầu vào
\newcommand{\bH}{\bm{H}}                   % Ma trận Hessian hoặc xấp xỉ BFGS
\newcommand{\bI}{\bm{I}}                   % Ma trận đơn vị
\newcommand{\bJ}{\bm{J}}                   % Ma trận Jacobi
\newcommand{\bK}{\bm{K}}                   % Kernel matrix
\newcommand{\bS}{\bm{S}}                   % Truy hồi Hessian theo biến đầu vào
\newcommand{\bU}{\bm{U}}
\newcommand{\bV}{\bm{V}}
\newcommand{\bW}{\bm{W}}                   % Ma trận trọng số (Weights)
\newcommand{\bX}{\bm{X}}                   % Ma trận dữ liệu
\newcommand{\bY}{\bm{Y}}
\newcommand{\bzero}{\bm{0}}                % Vector/Ma trận zero

% Ký hiệu Hy Lạp in đậm
\newcommand{\btheta}{\bm{\theta}}          % Tham số mạng
\newcommand{\bxi}{\bm{\xi}}                % Biến ngẫu nhiên / Điểm gieo
\newcommand{\bzeta}{\bm{\zeta}}
\newcommand{\bmu}{\bm{\mu}}
\newcommand{\bsigma}{\bm{\sigma}}
\newcommand{\blambda}{\bm{\lambda}}

\newcommand{\transpose}{^{\top}}           % Chuyển vị: \bx\transpose

%--------------------------------------------------------------
% 3. TOÁN TỬ GIẢI TÍCH & ĐẠI SỐ
%--------------------------------------------------------------
\newcommand{\dd}{\mathrm{d}}
\newcommand{\pd}[2]{\frac{\partial #1}{\partial #2}}
\newcommand{\pdd}[2]{\frac{\partial^{2} #1}{\partial #2^{2}}}
\newcommand{\od}[2]{\frac{\dd #1}{\dd #2}}
\newcommand{\odd}[2]{\frac{\dd^{2} #1}{\dd #2^{2}}}

\newcommand{\grad}{\nabla}
\newcommand{\gradt}{\nabla_{\btheta}}
\newcommand{\lap}{\Delta}

\DeclareMathOperator{\D}{D}                 % Ma trận Jacobi
\DeclareMathOperator{\diag}{diag}
\DeclareMathOperator{\tr}{tr}
\DeclareMathOperator{\rank}{rank}           % Hạng của ma trận
\DeclareMathOperator{\spanop}{span}         % Không gian sinh
\DeclareMathOperator{\supp}{supp}           % Giá (support) của hàm
\DeclareMathOperator{\Var}{Var}             % Phương sai
\DeclareMathOperator{\Cov}{Cov}             % Hiệp phương sai
\DeclareMathOperator*{\argmin}{arg\,min}
\DeclareMathOperator*{\argmax}{arg\,max}

% Epsilon máy (phân biệt với \varepsilon của thuật toán Adam)
\newcommand{\eps}{\varepsilon_{\mathrm{m}}}

%--------------------------------------------------------------
% 4. CHUẨN VÀ TRỊ TUYỆT ĐỐI (Yêu cầu package mathtools)
%--------------------------------------------------------------
\DeclarePairedDelimiter{\norm}{\lVert}{\rVert}
\DeclarePairedDelimiter{\inner}{\langle}{\rangle}

% Thêm trợ lý tự động thay đổi kích thước khi dùng dấu * (ví dụ \norm*{...})
% Hoặc dùng \norm[\big]{...}

%--------------------------------------------------------------
% 5. MẠNG NƠ-RON
%--------------------------------------------------------------
\DeclareMathOperator{\Net}{Net}             % \Net_{\btheta}(\bx)
\newcommand{\act}{\sigma}                   % Hàm kích hoạt
\newcommand{\layer}[1]{^{(#1)}}             % Tầng: \bW\layer{l}
\newcommand{\nlayer}{L}                     % Tổng số tầng
\newcommand{\lr}{\alpha}                    % Tốc độ học (Learning rate)

%--------------------------------------------------------------
% 6. HÀM MẤT MÁT VÀ CÁC THÀNH PHẦN CỦA PINN
%--------------------------------------------------------------
\newcommand{\Lop}{\mathcal{L}}             % Toán tử vi phân (viết tắt 1)
\newcommand{\diffop}{\mathcal{L}}          % Toán tử vi phân (viết tắt 2)
\newcommand{\bdryop}{\mathcal{B}}          % Toán tử điều kiện biên
\newcommand{\resid}{\mathcal{R}}           % Phần dư (Residual)
\newcommand{\ptloss}{\varrho}              % Mất mát tại một điểm

\newcommand{\Jr}{J_{\mathrm{r}}}            % Mất mát phần dư (PDE)
\newcommand{\Jb}{J_{\mathrm{b}}}            % Mất mát điều kiện biên
\newcommand{\Ji}{J_{\mathrm{ic}}}           % Mất mát điều kiện ban đầu
\newcommand{\Jd}{J_{\mathrm{data}}}         % Mất mát dữ liệu quan sát

% Trọng số các thành phần
\newcommand{\lamr}{\lambda_{\mathrm{r}}}
\newcommand{\lamb}{\lambda_{\mathrm{b}}}
\newcommand{\lami}{\lambda_{\mathrm{ic}}}
\newcommand{\lamd}{\lambda_{\mathrm{data}}}

% Tập điểm huấn luyện (Collocation points)
\newcommand{\setr}{\mathcal{T}_{\mathrm{r}}}
\newcommand{\setb}{\mathcal{T}_{\mathrm{b}}}
\newcommand{\seti}{\mathcal{T}_{\mathrm{ic}}}
\newcommand{\setd}{\mathcal{T}_{\mathrm{data}}}
\newcommand{\Nr}{N_{\mathrm{r}}}
\newcommand{\Nb}{N_{\mathrm{b}}}
\newcommand{\Ni}{N_{\mathrm{ic}}}
\newcommand{\Nd}{N_{\mathrm{data}}}

%--------------------------------------------------------------
% 7. SAI SỐ VÀ ĐÁNH GIÁ
%--------------------------------------------------------------
\newcommand{\Ltwoerr}{\epsilon_{L^{2}}}
\newcommand{\Linferr}{\epsilon_{L^{\infty}}}
\newcommand{\uex}{u^{\star}}                % Nghiệm chính xác (khớp ký hiệu Chương 5)
\newcommand{\uhat}{\hat{u}}                 % Nghiệm dự đoán

%---------------------------------------------------------------
\DeclareMathOperator{\sech}{sech}
\newcommand{\Bop}{\mathcal{B}}             % Toán tử biên (hoặc \mathbf{B} tùy bạn)
\newcommand{\Tset}{\mathcal{T}}            % Tập hợp điểm phối trí (Collocation points)